\section{引言}
\label{sec:intro}

部分子分布函数（PDF）在高能粒子物理中扮演着基础性角色，它编码了质子与中子内部夸克和胶子的动量分布。
这些函数是在诸如大型强子对撞机（LHC）等强子对撞机上作出精确预言的重要输入，直接影响从标准模型（SM）到超越标准模型（BSM）等一系列过程的诠释。
在各类 PDF 之中，胶子分布尤为关键，因为胶子在高能量下主导质子的动量，并且在希格斯玻色子产生与喷注形成等过程中起重要作用。
夸克 PDF 已被深度非弹性散射与 Drell--Yan 数据较好地约束，而胶子 PDF 的确定仍然不够精确，尤其是在中间及大动量分数（$x$）区域。
这一不确定性限制了 LHC 上关键观测量理论预言的精度。
因此，提高胶子 PDF 的精度有望显著增强 BSM 物理搜寻的灵敏度，并深化我们对强子对撞中量子色动力学（QCD）动力学的理解~\cite{Amoroso:2022eow}。

格点 QCD（LQCD）通过在有限格点上离散化时空，为研究强相互作用的非微扰区提供了第一性原理的方法。
尽管 LQCD 传统上仅限于计算少数强子结构观测量（例如 PDF 的矩），近年的突破已使 LQCD 计算能够触及 PDF 对 Bjorken-$x$ 的依赖关系。
大动量有效理论（LaMET）~\cite{Ji:2013dva} 与赝 PDF（pseudo-PDF）~\cite{Radyushkin:2017cyf} 是从格点可观测量中提取光锥关联函数最流行的两种方法。
过去十年左右，不仅在 PDF 领域，在其他种类的三维结构（如广义部分子分布 GPD）方面也取得了长足进展；
有兴趣的读者可参阅文献~\cite{Ji:2020ect,Constantinou:2020hdm,Lin:2023kxn}的综述。
然而，格点上依赖于 $x$ 的结果大多集中在同位矢量夸克分布上，关于胶子 PDF 的工作很少。

在格点 QCD 中提取胶子 PDF 仍然更具挑战性，原因在于其信号本质上噪声更大，且需要处理更复杂的算符混合。
尽管如此，近年来已有令人鼓舞的进展报道，大多采用赝 PDF 框架，得到了核子~\cite{Fan:2020cpa, Fan:2022kcb, HadStruc:2021wmh, Delmar:2023agv, HadStruc:2022yaw}、π介子~\cite{Good:2023ecp}与 K介子~\cite{Salas-Chavira:2021wui}的胶子 PDF。
然而，赝 PDF 方法往往需要借助唯象启发、依赖于 $x$ 的 PDF 模型来拟合格点矩阵元的比值；不过近期也有减少该方法模型依赖的努力~\cite{Karpie:2019eiq, Karpie:2021pap}。
相比之下，LaMET 先对格点矩阵元作傅里叶变换得到准 PDF（quasi-PDF），再经由匹配程序将其匹配到光锥 PDF，匹配精度按部分子动量的逆幂次收敛。
LaMET 中的模型依赖体现在为获得可靠的傅里叶变换而进行的大距离矩阵元外推上。
同样用 LaMET 得到胶子 PDF 将使学界极感兴趣，这样便可通过比较两种方法论的结果更好地控制格点系统误差。

最近，MSULat 组致力于用 LaMET 推进胶子 PDF 的 LQCD 计算~\cite{Good:2024iur}。
我们在格距约 0.12~fm 的系综（ensemble）上，于两个π介子质量（约 310 与 690~MeV）处进行了高统计量测量。
该工作探索了应用库仑规范固定来提升大 Wilson 线间距处的信号质量，展示了其提高胶子矩阵元计算可靠性的潜力。
通过分析多个胶子算符并比较核子态与π介子态所得矩阵元，该研究确定了提取胶子 PDF 最有效的算符。
然而，由于缺少对该算符进行混合比值重正化（hybrid-ratio renormalization）所需的微扰 QCD（pQCD）计算，我们当时只能以重正化的估计值进行分析。

在本工作中，我们取自前人工作~\cite{Good:2024iur}的格点胶子矩阵元，计算了混合重正化方案的匹配系数，从而用 LaMET 方法获得了首个光锥核子胶子 PDF。
我们在 Sec.~\ref{sec:analytic} 中给出胶子算符混合重正化方案的次领头阶（NLO）Wilson 系数。
我们在 Sec.~\ref{sec:numerical} 中确定混合重正化的抵消项，并将其应用于 LaMET 矩阵元的重正化。
我们将重正化矩阵元与唯象估计相比较，给出首个来自 LaMET 的核子胶子 PDF，并与唯象拟合进行比较。
我们在 Sec.~\ref{sec:conclusion} 中总结并展望胶子 PDF 的未来前景。
